% Generated draft: requires mathematical and source review.
\documentclass[11pt]{article}
\usepackage[margin=25mm]{geometry}
\usepackage{fontspec}
\setmainfont{Latin Modern Roman}
\usepackage{amsmath,amssymb,mathtools,unicode-math}
\setmathfont{Latin Modern Math}
\usepackage{enumitem,needspace,xurl,hyperref}
\setlength{\parindent}{0pt}
\setlength{\parskip}{5pt}
\setlength{\emergencystretch}{3em}
\newcommand{\N}{\mathbb{N}}
\newcommand{\Z}{\mathbb{Z}}
\newcommand{\Q}{\mathbb{Q}}
\newcommand{\R}{\mathbb{R}}
\newcommand{\C}{\mathbb{C}}
\newcommand{\F}{\mathbb{F}}
\newcommand{\A}{\mathbb{A}}
\newcommand{\PP}{\mathbb{P}}
\newcommand{\E}{\mathbb{E}}
\newcommand{\eps}{\varepsilon}
\newcommand{\dd}{\,\mathrm{d}}
\begin{document}
\section*{UnsolvedMath ID: 30000713}
\textbf{Source problem code:} catalogue entry
\subsection{Definitions and mathematical statement}
For each \(r\in\Z_{\geq0}\), equip
\[
 S^r=\{x\in\R^{r+1}:x_0^2+\cdots+x_r^2=1\}
\]
with its standard real algebraic structure and Euclidean topology. In particular \(S^0\) is the two-point sphere, and the integers \(n,p\) below include zero.

The source conjecture asks whether, for every continuous map \(f:S^n\to S^p\), there exist a regular map \(g:S^n\to S^p\) and a continuous homotopy
\[
 H:S^n\times[0,1]\longrightarrow S^p,
 \qquad H(x,0)=f(x),\quad H(x,1)=g(x).
\]
The quantifiers are \(\forall n,p\geq0\ \forall f\ \exists g,H\); the representative can depend on the map and dimension pair. No basepoint condition or regularity of the homotopy itself is required by the wording.

The term regular needs its real-algebraic convention. In a common real-algebraic category, it means coordinate functions locally given by rational functions with denominators nonzero on the relevant real neighborhoods. A polynomial restriction is regular, but an assertion specifically about polynomial representatives or morphisms of the full affine schemes can be stronger than this no-real-poles interpretation. The record does not define the convention, so these categories are not silently identified. The standard quadratic models above are likewise made explicit; replacing a sphere by a different algebraic model would be a different formulation without further comparison.

The task is existence of at least one regular representative in each continuous homotopy class. It is not a requirement that every continuous map already be regular, that regular maps be dense in all continuous maps, or that arbitrary regular representatives admit regular homotopies. Both \(n\) and \(p\) range over all nonnegative integers, without a dimension inequality.

The literature triage checked 2026-08-21 reports known results for selected dimension pairs and says the all-pairs assertion remains unresolved. It supplies no list or precise hypotheses of those cases. The open label and reported official-source verification remain dated metadata, not independently checked current conclusions.
\subsection{Short English statement}
Does every continuous homotopy class from any standard sphere to any other contain a regular real-algebraic map? The homotopy is only required to be continuous, and zero-dimensional spheres remain included. The source does not define whether regular uses the no-real-poles convention or the stronger polynomial/scheme setting.
\subsection{Sources}
\begin{itemize}
\item Local UnsolvedMath registry, permanent ID 30000713, source code catalogue entry, ``Regular Representatives of Homotopy Classes Between Spheres,'' including the original conjecture, all nonnegative dimension pairs and verification metadata.
\item \emph{Reelle Algebraische Geometrie (2007)}, Oberwolfach Reports, pp.\ 753--820, DOI 10.4171/owr/2007/14: \url{https://doi.org/10.4171/owr/2007/14}.
\item Supplied reference ``regular-map homotopy results'': \url{https://doi.org/10.1007/s00208-006-0061-3}.
\item Supplied reference ``later sphere-map advances'': \url{https://doi.org/10.1016/j.matpur.2022.03.002}.
\item Literature triage checked 2026-08-21; record updated 2026-09-16. No external source has been visited or independently checked here.
\end{itemize}
\end{document}
